\documentclass[12pt,a4paper]{article}
\usepackage{fontawesome5}
\providecommand{\faPencilAlt}{\faIcon{pencil-alt}}
\usepackage{styles/couleurs}
\usepackage{styles/mise_en_page}
\usepackage{styles/boites}
\usepackage{styles/maths}
\usepackage{styles/figures}
\usepackage{enumitem}
\usepackage{array}
\usepackage{tabularx}
\usepackage{amsmath}
\usepackage{amssymb}
\usepackage{tikz}
\usepackage{xparse}
\usetikzlibrary{calc,patterns,arrows.meta}
\setlist[itemize]{leftmargin=1.4em,itemsep=0.15em,topsep=0.2em}
\setlist[enumerate]{leftmargin=*,itemsep=0.25em,topsep=0.25em}
\newcommand{\lignepointillee}{\noindent\dotfill\par\vspace{0.45em}}
\newcommand{\lignesreponse}[1]{%
\begin{tcolorbox}[enhanced,colback=white,colframe=NeutreC!40,boxrule=0.5pt,
arc=1.5mm,left=3mm,right=3mm,top=2mm,bottom=1mm,before skip=2pt,after skip=2pt]
\foreach \i in {1,...,#1}{\lignepointillee}
\end{tcolorbox}}
\newcommand{\titreseq}[1]{%
\begin{center}
{\sffamily\large\bfseries\color{SavoirC}#1}
\end{center}}
\begin{document}

\titreseq{Activité 2 -- Le mur de briques}

Le contenu d'une brique est la somme des deux briques placées juste en dessous.
Les deux briques grises contiennent le même nombre, noté \(x\).

\begin{center}
\begin{tikzpicture}[scale=1.05]
  \def\w{1.7}
  \def\h{0.75}
  \foreach \i/\lab/\fillc in {0/2/white,1/\(x\)/NeutreC!20,2/\(x\)/NeutreC!20,3/3/white}{
    \fill[\fillc] (\i*\w,0) rectangle ++(\w,\h);
    \draw[line width=0.8pt] (\i*\w,0) rectangle ++(\w,\h);
    \node at (\i*\w+0.5*\w,0.5*\h) {\lab};
  }
  \foreach \i in {0,1,2}{\draw[line width=0.8pt] (\i*\w+0.5*\w,\h) rectangle ++(\w,\h);}
  \foreach \i in {0,1}{\draw[line width=0.8pt] (\i*\w+\w,2*\h) rectangle ++(\w,\h);}
  \draw[line width=0.8pt] (1.5*\w,3*\h) rectangle ++(\w,\h);
\end{tikzpicture}
\end{center}

\section*{Partie A -- Essayer avec des nombres}

Complète le mur pour \(x=1\) et \(x=10\).

\renewcommand{\arraystretch}{1.5}
\begin{tabularx}{0.8\textwidth}{|>{\centering\arraybackslash}X|>{\centering\arraybackslash}X|}
\hline
\textbf{Valeur de \(x\)} & \textbf{Sommet}\\
\hline
\(1\) & \\
\hline
\(10\) & \\
\hline
\end{tabularx}

\section*{Partie B -- Compléter avec des expressions}

Ecris une expression dans chaque brique.

\vspace{1cm}

\begin{center}
\begin{tikzpicture}[scale=2.5]
  \def\w{1.7}
  \def\h{0.75}
  \foreach \i/\lab/\fillc in {0/2/white,1/\(x\)/NeutreC!20,2/\(x\)/NeutreC!20,3/3/white}{
    \fill[\fillc] (\i*\w,0) rectangle ++(\w,\h);
    \draw[line width=0.8pt] (\i*\w,0) rectangle ++(\w,\h);
    \node at (\i*\w+0.5*\w,0.5*\h) {\lab};
  }
  \foreach \i in {0,1,2}{\draw[line width=0.8pt] (\i*\w+0.5*\w,\h) rectangle ++(\w,\h);}
  \foreach \i in {0,1}{\draw[line width=0.8pt] (\i*\w+\w,2*\h) rectangle ++(\w,\h);}
  \draw[line width=0.8pt] (1.5*\w,3*\h) rectangle ++(\w,\h);
\end{tikzpicture}
\end{center}

\section*{Partie C -- Trouver une formule courte}

Montre que le sommet peut s'écrire \(6x+5\).

Explique pourquoi \(x+x=2x\) et \(2x+2x+2x=6x\).

\lignesreponse{6}

\section*{Partie D -- Utiliser la formule}

Calcule rapidement le sommet pour \(x=100\), \(x=-15\) et \(x=2{,}5\).

\lignesreponse{3}

\end{document}
